% ---- pg-341 （书页 329） ----

\noindent (d) 对假设 $H_0:\beta_2=0$, 作显著性水平为 $\alpha=0.05$ 的似然比检验. 原来的模型是否应修改? 请讨论.

\medskip
\noindent 7.18\quad 计算 (a) 对练习 1.4 中的《福布斯》数据的各种可能的回归作 $C_p$ 图.

\noindent (b) 每种可能回归的 AIC.

\medskip
\noindent 7.19\quad 卫星应用推动了银-锌电池的发展. \hyperref[tab:7.5]{表 7.5} 中列出了表征电池在其寿命周期内的性能失效数据. 利用这些数据

\begin{mvtable}
\phantomsection\label{tab:7.5}\mvtabcap{表 7.5\quad 电池失效数据}
\footnotesize
\begin{tabular}{cccccc}
\toprule
\makecell{$Z_1$\\充电率\\（安）}&\makecell{$Z_2$\\放电率\\（安）}&
\makecell{$Z_3$\\放电深度\\（在安培小时计算\\值中所占百分比（\%））}&
\makecell{$Z_4$\\温度\\（℃）}&\makecell{$Z_5$\\充电电压极限\\（伏）}&
\makecell{$Y$\\ \\失效周期}\\
\midrule
0.375&3.13&60.0&40&2.00&101\\
1.000&3.13&76.8&30&1.99&141\\
1.000&3.13&60.0&20&2.00&96\\
1.000&3.13&60.0&20&1.98&125\\
1.625&3.13&43.2&10&2.01&43\\
1.625&3.13&60.0&20&2.00&16\\
1.625&3.13&60.0&20&2.02&188\\
0.375&5.00&76.8&10&2.01&10\\
1.000&5.00&43.2&10&1.99&3\\
1.000&5.00&43.2&30&2.01&386\\
1.000&5.00&100.0&20&2.00&45\\
1.625&5.00&76.8&10&1.99&2\\
0.375&1.25&76.8&10&2.01&76\\
1.000&1.25&43.2&10&1.99&78\\
1.000&1.25&76.8&30&2.00&160\\
1.000&1.25&60.0&0&2.00&3\\
1.625&1.25&43.2&30&1.99&216\\
1.625&1.25&60.0&20&2.00&73\\
0.375&3.13&76.8&30&1.99&314\\
0.375&3.13&60.0&20&2.00&170\\
\bottomrule
\end{tabular}
\end{mvtable}

\noindent{\footnotesize 资料来源: 选自 S. Sidik, H. Leibecki, J. Bozek, \textit{Failure of Silver-Zinc Cells with Competing Failure Modes-Preliminary Data Analysis}, NASA Technical Memorandum 81556 (Cleveland: Lewis Research Center, 1980).}

\medskip
\noindent (a) 求 $\ln(Y)$ 对一个适当的(“最佳的”)预测变量子集的线性回归方程(估计的).

\noindent (b) 对 (a) 中拟合的模型作残差图, 用以检查正态性假定.

\medskip
\noindent 7.20\quad 用\hyperref[tab:7.5]{表 7.5} 中的电池失效数据作 $\ln(Y)$ 对预测变量 $z_1,z_2,\cdots,z_5$ 的第一主成分的线性回归(见 8.3 节). 将结果与练习 7.19(a) 中拟合的模型作比较.

\medskip
\noindent 7.21\quad 考虑\hyperref[tab:1.5]{表 1.5} 中的空气污染数据. 令 $Y_1=\mathrm{NO}_2$, 和 $Y_2=\mathrm{O}_3$ 为两个响应变量(污染物), 相应预测变量 $Z_1=$ 风, 及 $Z_2=$ 太阳辐射.

\noindent (a) 只用第一个响应 $Y_1$ 作回归分析.

\noindent (i) 提出并拟合适当的线性回归模型.
